\documentclass[11pt,xcolor={dvipsnames},aspectratio=159,hyperref={pdftex,pdfpagemode=UseNone,hidelinks,pdfdisplaydoctitle=true},usepdftitle=false]{beamer}
\usepackage{presentation}[aspectratio=169]
\usepackage{math}
\usepackage{mathtools}
\usepackage{mleftright}
% Enter title of presentation PDF:
\hypersetup{pdftitle={Introduction}}
% Enter link to PDF file with figures:

\begin{document}
% Enter presentation title:
\title{Introduction}
\subtitle{Fiscal and Monetary Policy 2026}
% Enter presentation information:

% Enter presentation authors:
\author{Piotr Żoch}%
% Enter presentation location and date (optional; comment line if not needed):
\frame{\titlepage}

% Fill out content of presentation:


\begin{frame}{About this course}
    \begin{itemize}
    \item \al{Advanced graduate course} on fiscal and monetary policy. 
    \item Not an exhaustive survey of the literature, but rather a selection of topics that I find interesting and important.
    \item The course will be based on a mix of theoretical and empirical papers.
    \item We will cover both tools and substance.
    \item Prerequisite: I will assume you took a course in Advanced Macroeconomics and are familiar with basic tools of dynamic macroeconomic analysis.
    \item Goal: you should be able to read and understand the most recent research in the field. 
    \end{itemize}
\end{frame}

\begin{frame}{Logistics}
    \begin{itemize}
    \item My email address: \url{p.zoch@uw.edu.pl}
    \item Office hours: by appointment
    \item All materials will be posted on the course website: \url{https://github.com/pzoch/FMP2026}
    \end{itemize}
\end{frame}


\begin{frame}{Roadmap}
    \begin{enumerate}
    \item Government spending multiplier: can government purchases stimulate the economy? 
    \item Fiscal capacity: what determines the ability of the government to raise revenue by issuing debt?
    \item Debt management: how should deficits and debt respond to macroeconomic shocks?
    \item Price level determination in monetary economies: what pins down the price level and determines the inflation rate?
    \item Fiscal and monetary policy interactions: how actions of the monetary authority affect the fiscal authority, and vice versa?
    \end{enumerate}
\end{frame}

\begin{frame}{Requirements}
    \begin{enumerate}
    \item \al{Problem sets}: two or three problem sets that will ask you to derive conclusions from the models we discuss in class. You can work in groups of up to two students.
    \item \alg{Final exam}: there will be a final exam in the examination session.
    \end{enumerate}
\end{frame}


\end{document}